\section{Background}
\label{sec:background}

\subsection{Parallel Tempering in Statistical Mechanics}

Parallel Tempering (PT) was introduced by \citet{swendsen1986} for simulating
spin-glass systems at low temperature.
At low temperature, a single Markov chain may become trapped in a local energy
minimum for exponentially long times, exploring only a small fraction of the
configuration space.
Running multiple chains at increasing temperatures simultaneously allows high-temperature
chains to mix rapidly and periodically propose their configurations as candidates
for the low-temperature chain.

\citet{geyer1991} formalised PT as a statistical Monte Carlo method,
establishing that the joint chain over all replicas is ergodic and that the
marginal distribution of the cold replica is exactly the target distribution.
This is the theoretical foundation for using only the cold chain's output in
plan selection: the hot chains are sampling aids, not target distributions.

\paragraph{Temperature and mixing.}
In a physical system with energy $E(\sigma)$ and inverse temperature $\beta$,
the Boltzmann acceptance criterion is $\min(1, \exp(-\beta \Delta E))$.
High temperature (low $\beta$) means almost all proposals are accepted;
the chain moves freely.
Low temperature (high $\beta$) means only improvements or small increases are
accepted; the chain explores a narrow basin of the energy landscape.

The replica exchange criterion between replica $i$ (at $\beta_i$) and replica
$j$ (at $\beta_j > \beta_i$, so $\varepsilon_j > \varepsilon_i$) is:
\[
  P(\text{swap}) = \min\!\bigl(1,\, \exp\bigl((\beta_i - \beta_j)(E_i - E_j)\bigr)\bigr),
\]
where $E_i = \EC(\text{assignment}_i)$ is the edge cut of replica $i$'s current
plan.
This is a valid Metropolis proposal on the joint distribution over all replicas,
preserving the product-of-marginals target distribution~\citep{geyer1991}.

\paragraph{Balance tolerance as redistricting temperature.}
In redistricting, the population balance tolerance $\varepsilon$ plays the role
of temperature.
A plan assignment $\mathbf{a}: V \to \{1,\ldots,k\}$ is \emph{feasible at
tolerance $\varepsilon$} if every district satisfies:
\[
  \left|\frac{\text{pop}(d)}{\bar{p}} - 1\right| \le \varepsilon,
\]
where $\bar{p} = P_{\mathrm{total}}/k$ is the ideal population per district.
At $\varepsilon = 0.005$ (statutory cold tolerance), only plans with $\pm 0.5\%$
population deviation are accessible.
At $\varepsilon = 0.05$ (hot tolerance), plans with up to $\pm 5\%$ deviation
are accessible---a much larger feasible set.
The \fr{} chain at high tolerance can make proposals that jump between
topologically distinct plan types (e.g., moving a district's centre of gravity
by several tracts in a single spanning-tree resampling), while the cold chain
can only make local adjustments within the tight balance constraint.

\paragraph{The ${\approx}23\%$ swap acceptance target.}
\citet{kone2005} showed that for Gaussian continuous distributions, the
optimal swap acceptance rate between adjacent replicas is approximately $23\%$,
achieved by spacing temperatures so that $\Delta(\ln T) \approx$ const.
For discrete combinatorial chains, this result is an empirical guideline rather
than a derived theorem.
Nevertheless, we use it as a calibration target for the geometric tolerance
ladder: if observed swap acceptance rates are near $23\%$ across all adjacent
pairs, the ladder is well-spaced.
Section~\ref{sec:empirical} reports observed rates of $18$--$24\%$ across NC,
WI, and TX, consistent with this target.

\subsection{ForestRecom as the Per-Replica Chain}

Each replica in \pt{} redistricting runs one instance of \fr{}
\citep{dellaLibera2026forestRecom}, the spanning-forest recombination chain.
\fr{} proposes plan changes by:
\begin{enumerate}
\item Selecting two adjacent districts $d_i, d_j$.
\item Sampling a random spanning forest of the induced subgraph $G[d_i \cup d_j]$
  with one component per district (a forest with exactly one edge cut between
  the two components).
\item Accepting the induced partition if the population balance constraint is
  satisfied at the replica's tolerance.
\end{enumerate}
\fr{} targets an approximately uniform distribution over balanced plans at its
configured tolerance~\citep{dellaLibera2026forestRecom}.
Each replica in \pt{} targets this uniform distribution at its own tolerance
level; the cold chain's stationary distribution is the target we care about.

\paragraph{Forward and reverse RNG split.}
\fr{} uses separate forward and reverse random number streams (for the forward
proposal and the reverse-proposal computation in the Hastings ratio).
In \pt{}, each replica's forward and reverse streams are derived from the
replica seed via SHA-256 with distinct prefixes, as specified in
Section~\ref{sec:seeding}.

\subsection{Search-Layer Context in the B-Series Platform}

The B-series redistricting platform uses a three-layer compositor:
\emph{Structure} (how the bisection tree is built),
\emph{WeightsOverride} (what edge weights signal),
and \emph{SeedCompositor} / \emph{Search} (how the plan space is explored and
the best plan selected).
\pt{} is a Search-layer algorithm.

\paragraph{Existing search strategies.}
\texttt{single}: Run one \fr{} chain for $T$ steps, return the minimum-EC plan.
\texttt{multi}: Run $T$ independent chains from different seeds, return the
minimum-EC plan overall.
\texttt{convergence} (\cs{}, B.16): Run chains until $T=600$ consecutive seeds
produce no improvement.
\texttt{percentile}: Run one chain, return the plan at the $p$-th percentile
of the edge-cut distribution.
\texttt{bisection-ensemble}: Run chains over a bisection-tree ensemble.

\pt{} is a new Search strategy: run $N$ chains at different tolerances
simultaneously, with periodic swap proposals, and return the plan at the
$p$-th percentile of the \emph{cold chain's} edge-cut distribution.

\paragraph{Orthogonality.}
\pt{} is orthogonal to the Structure and WeightsOverride layers.
It can be combined with any structure configuration (\geosec{},
\ar{}, standard bisect) and any weight configuration (geographic,
county, unweighted).
The empirical experiments in Section~\ref{sec:empirical} use the
\geosec{} structure and geographic weights as the standard baseline.

\subsection{Comparison with Multi-Scale ReCom}

\ms{}~\citep{dellaLibera2026multiscale} and \pt{} are both designed to improve
mixing for large-$k$ states.
Their mechanisms differ:

\begin{center}
\renewcommand{\arraystretch}{1.25}
\begin{tabular}{@{}lll@{}}
\toprule
Property & \ms{} (G.11) & \pt{} (B.20) \\
\midrule
Mixing mechanism & Geographic coarsening & Temperature ladders \\
Proposal size & Large (super-tract level) & Standard (tract level) \\
Approximation & Yes (coarsening resolution) & No (exact \fr{} at each $\varepsilon$) \\
Chain count & 1 & $N_{\mathrm{replicas}}$ \\
Extra cost & $\alpha \times$ coarsening overhead & $N_{\mathrm{replicas}} \times$ per-step cost \\
Best for & Strong geographic structure & Tolerance-induced barriers \\
\bottomrule
\end{tabular}
\end{center}

The two approaches are complementary and can in principle be combined (each PT
replica running an MS chain at its tolerance), though this is left to future
work.
In practice, PT is preferred when the bottleneck is the balance-tolerance
constraint rather than geographic coarsening, and MS is preferred when the
bottleneck is proposal size.
